\chapter{图表与代码}\label{chap:float}

本章介绍插图、表格、算法和代码等浮动体的使用方法。浮动体可以自由浮动在页面中，\LaTeX{} 会自动选择合适的位置进行排版。如需固定位置，可使用 \pkg{float} 宏包的 \texttt{[H]} 参数。

浮动体位置选项说明：\texttt{h}（here，当前位置）、\texttt{t}（top，页面顶部）、\texttt{b}（bottom，页面底部）、\texttt{p}（page，独立浮动页）、\texttt{!}（放宽内部限制）、\texttt{H}（强制固定位置，需 \pkg{float} 宏包）。

\section{插图}

\subsection{单个图形}

要插入单个图形，可以使用 \verb|\includegraphics| 命令，并在其参数中指定图像文件名和一些选项。例如：

\begin{figure}[!htb]
  \centering
  \includegraphics[width=0.7\textwidth]{figures/tongjithesis.pdf}
  \caption{\TongjiThesis{} 项目标志}
  \label{fig:example}
\end{figure}

这段代码会在浮动体中插入一个名为 \verb|figures/tongjithesis.pdf| 的图形，并将其居中显示。图形下方会有一个标题，同时为该图形标上标签，以便于交叉引用。

如需在正文中引用该图，可使用 \verb|\cref{fig:example}| 命令，效果为“\cref{fig:example}”；也可使用 \verb|\figref{fig:example}| 命令，效果为“\figref{fig:example}”。

需要注意的是，在使用 \verb|\includegraphics| 命令时，应该将图像文件放在与主文档相同的目录中，或者在参数中指定文件的完整路径。

推荐使用矢量图格式（PDF）用于图表和流程图，位图建议使用 PNG 格式（照片可用 JPG），矢量图在缩放时不会失真。

注意在浮动体内部应使用 \verb|\centering| 命令而非 \verb|\begin{center}...\end{center}| 环境，后者会引入额外的垂直间距。

\subsection{多个图形}

如果需要在同一浮动体中插入多个图形，可以使用 \pkg{subcaption} 宏包（旧的 \pkg{subfigure} 与 \pkg{subfig} 宏包已过时，请勿使用）。
最简单的做法是将多张图并排放入同一 \texttt{figure} 环境、用 \cs{hspace} 间隔——此时多图共用一个图号，没有各自子图题，如\cref{fig:example-subfig-1}。

\begin{figure}[!htb]
  \centering
  \includegraphics[width=0.4\textwidth]{example-image-a}
  \hspace{1cm}
  \includegraphics[width=0.4\textwidth]{example-image-b}
  \caption{这是一个示例图形组}
  \label{fig:example-subfig-1}
\end{figure}

如果多个图形相互独立，不共用一个图形计数器，可以使用 \texttt{minipage} 或 \texttt{parbox}，如\cref{fig:parallel1,fig:parallel2}。

\begin{figure}[!htb]
  \begin{minipage}{0.4\textwidth}
    \centering
    \includegraphics[width=\textwidth]{example-image-a}
    \caption{并排第一个图}
    \label{fig:parallel1}
  \end{minipage}\hfill
  \begin{minipage}{0.4\textwidth}
    \centering
    \includegraphics[width=\textwidth]{example-image-b}
    \caption{并排第二个图}
    \label{fig:parallel2}
  \end{minipage}
\end{figure}

若希望多个子图共用一个计数器并分别拥有子图题，使用 \pkg{subcaption} 宏包的 \cs{subcaptionbox} 命令（或 \texttt{subfigure} 环境），子图号以 (a)、(b) 区分，如\cref{fig:subcaptionbox}。

\begin{figure}[!htb]
  \centering
  \subcaptionbox{左图}{\includegraphics[width=0.4\textwidth]{example-image-a}}
  \hspace{1cm}
  \subcaptionbox{右图}{\includegraphics[width=0.3\textwidth]{example-image-b}}
  \caption{包含子图题的示例（使用 subcaptionbox）}
  \label{fig:subcaptionbox}
\end{figure}


\section{表格}

\subsection{基本表格}

编排表格时应简单明了、表达一致，内容明晰易懂，表文呼应，内容一致。表题应置于表格上方。

表格的编排建议采用国际通行的三线表\footnote{三线表以其形式简洁、功能分明、阅读方便而在科技论文中被推荐使用。三线表通常只有3条线，即顶线、底线和栏目线，没有竖线。}。可以使用 \pkg{booktabs} 宏包提供的 \cs{toprule}、\cs{midrule} 和 \cs{bottomrule} 命令来绘制三线表。同时，\pkg{longtable} 宏包可以与三线表很好地配合使用，可以排版较长的表格。
要编写三线表，需要在表格的头部、底部和每一栏之间绘制水平线。其中，\cs{toprule} 绘制顶部线，\cs{midrule} 绘制中部线，\cs{bottomrule} 绘制底部线。可以按照如下方式绘制三线表：

\begin{table}[!hpt]
  \caption{一个三线表}
  \label{tab:firstone}
  \centering
  \begin{tabular}{@{}llr@{}} \toprule
    \multicolumn{2}{c}{\textbf{Item}}                            \\ \cmidrule(r){1-2}
    \textbf{Animal} & \textbf{Description} & \textbf{Price} (\$) \\ \midrule
    Gnat            & per gram             & 13.65               \\
                    & each                 & 0.01                \\
    Gnu             & stuffed              & 92.50               \\
    Emu             & stuffed              & 33.33               \\
    Armadillo       & frozen               & 8.99                \\ \bottomrule
  \end{tabular}
\end{table}

\subsection{复杂表格}

我们经常会在表格下方标注数据来源，或者对表格里面的条目进行解释。可以用
\pkg{threeparttable} 实现带有脚注的表格，如\cref{tab:bigtable}。

\begin{table}[!htb]
  \centering
  \begin{threeparttable}[b]
    \caption{一个更大的三线表}
    \label{tab:bigtable}
    \begin{tabular}{@{}llrrrrrrr@{}} \toprule
      \multicolumn{2}{c}{\textbf{Item}} & \multicolumn{2}{c}{\textbf{Category 1}} & \multicolumn{2}{c}{\textbf{Category 2}} & \multicolumn{2}{c}{\textbf{Category 3}} & \multicolumn{1}{c}{\textbf{Total}}                                                                                     \\ \cmidrule(r){1-2} \cmidrule(lr){3-4} \cmidrule(lr){5-6} \cmidrule(lr){7-8} \cmidrule(l){9-9}
      \textbf{Animal}                   & \textbf{Description}                    & \textbf{Price} (\$)                     & \textbf{Quantity}                       & \textbf{Price} (\$)                & \textbf{Quantity} & \textbf{Price} (\$) & \textbf{Quantity} & \textbf{Price} (\$) \\ \midrule
      Gnat                              & per gram\tnote{a}                       & 13.65                                   & 100                                     & 12.35                              & 200               & 11.55               & 150               & 3650.00             \\
                                        & each                                    & 0.01                                    & 5000                                    & 0.01                               & 10000             & 0.009               & 20000             & 550.00              \\
      Gnu                               & stuffed                                 & 92.50                                   & 10                                      & 94.50                              & 15                & 96.50               & 20                & 5815.00             \\
      Emu                               & stuffed                                 & 33.33                                   & 25                                      & 34.33                              & 30                & 35.33               & 35                & 2704.95             \\
      Armadillo                         & frozen                                  & 8.99                                    & 50                                      & 7.99                               & 40                & 6.99                & 30\tnote{b}       & 1094.50             \\ \bottomrule
    \end{tabular}
    \begin{tablenotes}
      \item [a] 第一条脚注。
      \item [b] 第二条脚注。
    \end{tablenotes}
  \end{threeparttable}
\end{table}

为了编写更复杂的表格，我们可以使用 \href{https://www.tablesgenerator.com/}{Tables Generator} 来生成 \LaTeX{} 代码。该工具可以方便地设置表格的列数、行数、内容、格式、颜色等属性，支持常用的表格类型和样式，并且可以实时生成 \LaTeX{} 代码和预览效果，避免手动编写 \LaTeX{} 代码时出错。


如某个表需要转页接排，可以用 \pkg{longtable} 实现。接排时表题省略，表头应重复书
写，并在右上方写“续表 X”，如\cref{tab:performance}。

\begin{ThreePartTable}
  \begin{TableNotes}
    \item[a] 一个脚注
    \item[b] 另一个脚注
  \end{TableNotes}
  \begin{longtable}[c]{c*{6}{r}}
    \caption{实验数据}
    \label{tab:performance}                                                                                                                      \\
    \toprule
    \textbf{测试程序\tnote{b}} & \multicolumn{1}{c}{\textbf{正常运行}}   & \multicolumn{1}{c}{\textbf{同步}}
                           & \multicolumn{1}{c}{\textbf{检查点}}    & \multicolumn{1}{c}{\textbf{卷回恢复}}
                           & \multicolumn{1}{c}{\textbf{进程迁移}}   & \multicolumn{1}{c}{\textbf{检查点}}                                              \\
                           & \multicolumn{1}{c}{\textbf{时间} (s)} & \multicolumn{1}{c}{\textbf{时间} (s)}
                           & \multicolumn{1}{c}{\textbf{时间} (s)} & \multicolumn{1}{c}{\textbf{时间} (s)}
                           & \multicolumn{1}{c}{\textbf{时间} (s)} & \textbf{文件}（KB）                                                               \\
    \midrule
    \endfirsthead
    \tjlongtablecont{7}

    \toprule
    \textbf{测试程序\tnote{b}} & \multicolumn{1}{c}{\textbf{正常运行}}   & \multicolumn{1}{c}{\textbf{同步}}
                           & \multicolumn{1}{c}{\textbf{检查点}}    & \multicolumn{1}{c}{\textbf{卷回恢复}}
                           & \multicolumn{1}{c}{\textbf{进程迁移}}   & \multicolumn{1}{c}{\textbf{检查点}}                                              \\
                           & \multicolumn{1}{c}{\textbf{时间} (s)} & \multicolumn{1}{c}{\textbf{时间} (s)}
                           & \multicolumn{1}{c}{\textbf{时间} (s)} & \multicolumn{1}{c}{\textbf{时间} (s)}
                           & \multicolumn{1}{c}{\textbf{时间} (s)} & \textbf{文件}（KB）                                                               \\
    \midrule
    \endhead
    \hline
    \multicolumn{7}{r}{续下页}
    \endfoot
    \insertTableNotes
    \endlastfoot
    CG.C.2                 & 23.05                               & 0.002                               & 0.116          & 0.035 & 0.589 & 32491  \\
    CG.A.4                 & 15.06                               & 0.003                               & 0.067          & 0.021 & 0.351 & 18211  \\
    CG.A.8                 & 13.38                               & 0.004                               & 0.072          & 0.023 & 0.210 & 9890   \\
    CG.B.2                 & 867.45                              & 0.002                               & 0.864          & 0.232 & 3.256 & 228562 \\
    CG.B.4                 & 501.61                              & 0.003                               & 0.438          & 0.136 & 2.075 & 123862 \\
    CG.B.8                 & 384.65                              & 0.004                               & 0.457          & 0.108 & 1.235 & 63777  \\
    MG.A.2                 & 112.27                              & 0.002                               & 0.846\tnote{a} & 0.237 & 3.930 & 236473 \\
    MG.A.4                 & 59.84                               & 0.003                               & 0.442          & 0.128 & 2.070 & 123875 \\
    MG.A.8                 & 31.38                               & 0.003                               & 0.476          & 0.114 & 1.041 & 60627  \\
    MG.B.2                 & 526.28                              & 0.002                               & 0.821          & 0.238 & 4.176 & 236635 \\
    MG.B.4                 & 280.11                              & 0.003                               & 0.432          & 0.130 & 1.706 & 123793 \\
    MG.B.8                 & 148.29                              & 0.003                               & 0.442          & 0.116 & 0.893 & 60600  \\
    LU.A.2                 & 2116.54                             & 0.002                               & 0.110          & 0.030 & 0.532 & 28754  \\
    LU.A.4                 & 1102.50                             & 0.002                               & 0.069          & 0.017 & 0.255 & 14915  \\
    LU.A.8                 & 574.47                              & 0.003                               & 0.067          & 0.016 & 0.192 & 8655   \\
    LU.B.2                 & 9712.87                             & 0.002                               & 0.357          & 0.104 & 1.734 & 101975 \\
    LU.B.4                 & 4757.80                             & 0.003                               & 0.190          & 0.056 & 0.808 & 53522  \\
    LU.B.8                 & 2444.05                             & 0.004                               & 0.222          & 0.057 & 0.548 & 30134  \\
    EP.A.2                 & 123.81                              & 0.002                               & 0.010          & 0.003 & 0.074 & 1834   \\
    EP.A.4                 & 61.92                               & 0.003                               & 0.011          & 0.004 & 0.073 & 1743   \\
    EP.A.8                 & 31.06                               & 0.004                               & 0.017          & 0.005 & 0.073 & 1661   \\
    EP.B.2                 & 495.49                              & 0.001                               & 0.009          & 0.003 & 0.196 & 2011   \\
    SP.B.4                 & 397.69                              & 0.002                               & 0.015          & 0.005 & 0.122 & 1763   \\
    SP.B.8                 & 196.74                              & 0.003                               & 0.018          & 0.006 & 0.082 & 1865   \\
    AA.A.2                 & 13.81                               & 0.002                               & 0.012          & 0.004 & 0.074 & 1362   \\
    AA.A.4                 & 6.92                                & 0.003                               & 0.011          & 0.003 & 0.073 & 1331   \\
    AA.A.8                 & 3.06                                & 0.004                               & 0.017          & 0.004 & 0.073 & 1225   \\
    AA.B.2                 & 49.49                               & 0.001                               & 0.009          & 0.003 & 0.196 & 2254   \\
    AA.B.4                 & 24.69                               & 0.002                               & 0.012          & 0.004 & 0.122 & 1453   \\
    AA.B.8                 & 12.74                               & 0.003                               & 0.018          & 0.005 & 0.082 & 1432   \\
    AA.C.2                 & 13.81                               & 0.002                               & 0.012          & 0.004 & 0.074 & 1362   \\
    AA.C.4                 & 6.92                                & 0.003                               & 0.011          & 0.003 & 0.073 & 1331   \\
    AA.C.8                 & 3.06                                & 0.004                               & 0.017          & 0.004 & 0.073 & 1225   \\
    EP.C.2                 & 495.49                              & 0.001                               & 0.009          & 0.003 & 0.196 & 2011   \\
    EP.C.4                 & 397.69                              & 0.002                               & 0.015          & 0.005 & 0.122 & 1763   \\
    EP.C.8                 & 196.74                              & 0.003                               & 0.018          & 0.006 & 0.082 & 1865   \\
    \bottomrule
  \end{longtable}
\end{ThreePartTable}



\section{算法}

本模板默认使用 \pkg{algpseudocode} 排版算法，也可通过文档类选项 \texttt{algo=algorithm2e} 切换为 \pkg{algorithm2e}。

\iftjalgorithmtwoe
  \cref{algo:algorithm} 是一个使用 \pkg{algorithm2e} 的示例。

  \SetKw{Return}{return}
  \begin{algorithm}[H]
    \caption{计算斐波那契数列}\label{algo:algorithm}
    \SetKwFunction{Fib}{Fibonacci}
    \KwIn{$n \geq 0$}
    \KwOut{$fib(n)$}
    \Fib{$n$}\;
    \If{$n \leq 1$}{\Return $n$\;}
    \Return \Fib{$n-1$} + \Fib{$n-2$}\;
  \end{algorithm}
\else
  \cref{algo:algorithm} 是一个使用 \pkg{algpseudocode} 的示例。

  \begin{algorithm}
    \caption{计算斐波那契数列}\label{algo:algorithm}
    \begin{algorithmic}[1]
      \Require $n \geq 0$
      \Ensure $fib(n)$
      \Function{Fibonacci}{$n$}
      \If{$n \leq 1$}
      \State \Return $n$
      \Else
      \State \Return \Call{Fibonacci}{$n-1$} + \Call{Fibonacci}{$n-2$}
      \EndIf
      \EndFunction
    \end{algorithmic}
  \end{algorithm}
\fi

无论使用 \pkg{algpseudocode} 还是 \pkg{algorithm2e}，都可以通过相应的命令来定义函数、条件语句和循环结构，以便于编写清晰易懂的算法描述。

\section{代码}

\subsection{代码高亮方式}

本模板提供 \pkg{listings} 和 \pkg{minted} 两种代码高亮方案，可通过文档类选项 \texttt{minted=false}（默认）或 \texttt{minted=true} 切换。\pkg{listings} 是 \LaTeX{} 原生的代码显示包，始终可用，无需外部依赖；\pkg{minted} 基于 Python Pygments，提供更精确的语法解析和更丰富的高亮效果，但需要安装 Python 3.11–3.13 和 Pygments（\texttt{pip install pygments}），编译时需添加 \texttt{-shell-escape} 参数。若需要更好的高亮效果，可设置 \texttt{minted=true} 启用 \pkg{minted}。

\subsection{代码块示例}

\pkg{listings} 包是 \LaTeX{} 最常用的代码高亮工具，无论 \texttt{minted} 选项如何设置，它始终可用。如\cref{lst:basic-example} 所示。

\begin{lstlisting}[
  language=Python,         % 指定语法高亮的编程语言
  caption={使用listings包的Python示例},
  label=lst:basic-example,
  float
]
# 这是一个简单的Python函数
def hello_world(name: str = "World") -> str:
    """返回个性化问候语"""
    greeting = f"Hello, {name}!"
    return greeting

# 调用函数并打印结果
print(hello_world())  # 输出: Hello, World!
print(hello_world("LaTeX"))  # 输出: Hello, LaTeX!
\end{lstlisting}

\pkg{listings} 包支持绝大多数编程语言，如 C/C++、Java、Python、MATLAB、R、SQL 等，使用 \texttt{language} 参数指定即可。

\subsubsection{关于跨页代码段的说明}
\pkg{listings} 和 \pkg{minted} 均支持跨页显示代码段。一般而言，学术写作中不建议在正文中放置大段代码，建议仅展示关键代码片段，将完整代码放入附录或以外部链接形式提供。

本模板提供了 \texttt{longlisting} 环境，用于排版需要跨页显示并带有标题的代码块。与浮动体代码环境不同，\texttt{longlisting} 不会浮动，可以自然跨越多页。在环境内部放置代码块，并用 \cs{caption} 添加标题即可。以下为示例：

\makeatletter
\begin{longlisting}
\iftongjithesis@minted
\begin{minted}{python}
from dataclasses import dataclass, field
from typing import Dict, List, Tuple

# 同济大学五级制绩点对照（2017年版）
GRADE_TABLE: Dict[str, Tuple[int, int]] = {
    "A": (5, 95),  # 优 Excellent
    "B": (4, 85),  # 良 Good
    "C": (3, 75),  # 中 Fair
    "D": (2, 65),  # 及格 Pass
    "F": (0, 30),  # 不及格 Failed
}

VALID_GRADES = set(GRADE_TABLE.keys())


@dataclass
class Course:
    """课程信息"""
    name: str      # 课程名称
    credits: float # 学分
    grade: str     # 成绩等级 (A/B/C/D/F)

    def __post_init__(self):
        if self.grade not in VALID_GRADES:
            raise ValueError(f"无效的成绩等级: {self.grade}")
        if self.credits <= 0:
            raise ValueError(f"学分须为正数: {self.credits}")

    @property
    def grade_point(self) -> int:
        """课程绩点"""
        return GRADE_TABLE[self.grade][0]

    @property
    def equiv_score(self) -> int:
        """百分制折算值"""
        return GRADE_TABLE[self.grade][1]

    def __repr__(self) -> str:
        return f"{self.name}({self.credits}学分, {self.grade})"


@dataclass
class Transcript:
    """成绩单：按同济大学本科生成绩计算方法"""
    student_name: str = ""
    student_id: str = ""
    courses: List[Course] = field(default_factory=list)

    def add(self, name: str, credits: float, grade: str) -> None:
        """添加一门课程成绩"""
        self.courses.append(Course(name, credits, grade))

    @property
    def total_credits(self) -> float:
        """总学分"""
        return sum(c.credits for c in self.courses)

    @property
    def gpa(self) -> float:
        """平均绩点 = sum(学分 * 绩点) / sum(学分)"""
        if self.total_credits == 0:
            return 0.0
        weighted = sum(c.credits * c.grade_point for c in self.courses)
        return round(weighted / self.total_credits, 2)

    @property
    def avg_score(self) -> float:
        """百分制平均成绩 = sum(折算值 * 学分) / sum(学分)"""
        if self.total_credits == 0:
            return 0.0
        weighted = sum(c.credits * c.equiv_score for c in self.courses)
        return round(weighted / self.total_credits, 1)

    def ranking(self) -> List[Tuple[str, float, int]]:
        """按绩点降序排列课程"""
        return sorted(
            [(c.name, c.credits, c.grade_point) for c in self.courses],
            key=lambda x: x[2], reverse=True,
        )

    def summary(self) -> str:
        """生成成绩摘要"""
        lines = [f"学生: {self.student_name} ({self.student_id})"]
        lines.append(f"总学分: {self.total_credits}")
        lines.append(f"平均绩点 (GPA): {self.gpa}")
        lines.append(f"百分制平均成绩: {self.avg_score}")
        return "\n".join(lines)


if __name__ == "__main__":
    t = Transcript("张三", "2024001")
    t.add("高等数学", 6.0, "A")
    t.add("线性代数", 4.0, "B")
    t.add("数据结构", 4.0, "A")
    t.add("大学物理", 4.0, "C")
    t.add("程序设计", 3.0, "A")

    print(t.summary())
    print("\n课程排名:")
    for name, credits, gp in t.ranking():
        print(f"  {name}: {credits}学分, 绩点{gp}")
\end{minted}
\else
\begin{lstlisting}[language=Python]
from dataclasses import dataclass, field
from typing import Dict, List, Tuple

# 同济大学五级制绩点对照（2017年版）
GRADE_TABLE: Dict[str, Tuple[int, int]] = {
    "A": (5, 95),  # 优 Excellent
    "B": (4, 85),  # 良 Good
    "C": (3, 75),  # 中 Fair
    "D": (2, 65),  # 及格 Pass
    "F": (0, 30),  # 不及格 Failed
}

VALID_GRADES = set(GRADE_TABLE.keys())


@dataclass
class Course:
    """课程信息"""
    name: str      # 课程名称
    credits: float # 学分
    grade: str     # 成绩等级 (A/B/C/D/F)

    def __post_init__(self):
        if self.grade not in VALID_GRADES:
            raise ValueError(f"无效的成绩等级: {self.grade}")
        if self.credits <= 0:
            raise ValueError(f"学分须为正数: {self.credits}")

    @property
    def grade_point(self) -> int:
        """课程绩点"""
        return GRADE_TABLE[self.grade][0]

    @property
    def equiv_score(self) -> int:
        """百分制折算值"""
        return GRADE_TABLE[self.grade][1]

    def __repr__(self) -> str:
        return f"{self.name}({self.credits}学分, {self.grade})"


@dataclass
class Transcript:
    """成绩单：按同济大学本科生成绩计算方法"""
    student_name: str = ""
    student_id: str = ""
    courses: List[Course] = field(default_factory=list)

    def add(self, name: str, credits: float, grade: str) -> None:
        """添加一门课程成绩"""
        self.courses.append(Course(name, credits, grade))

    @property
    def total_credits(self) -> float:
        """总学分"""
        return sum(c.credits for c in self.courses)

    @property
    def gpa(self) -> float:
        """平均绩点 = sum(学分 * 绩点) / sum(学分)"""
        if self.total_credits == 0:
            return 0.0
        weighted = sum(c.credits * c.grade_point for c in self.courses)
        return round(weighted / self.total_credits, 2)

    @property
    def avg_score(self) -> float:
        """百分制平均成绩 = sum(折算值 * 学分) / sum(学分)"""
        if self.total_credits == 0:
            return 0.0
        weighted = sum(c.credits * c.equiv_score for c in self.courses)
        return round(weighted / self.total_credits, 1)

    def ranking(self) -> List[Tuple[str, float, int]]:
        """按绩点降序排列课程"""
        return sorted(
            [(c.name, c.credits, c.grade_point) for c in self.courses],
            key=lambda x: x[2], reverse=True,
        )

    def summary(self) -> str:
        """生成成绩摘要"""
        lines = [f"学生: {self.student_name} ({self.student_id})"]
        lines.append(f"总学分: {self.total_credits}")
        lines.append(f"平均绩点 (GPA): {self.gpa}")
        lines.append(f"百分制平均成绩: {self.avg_score}")
        return "\n".join(lines)


if __name__ == "__main__":
    t = Transcript("张三", "2024001")
    t.add("高等数学", 6.0, "A")
    t.add("线性代数", 4.0, "B")
    t.add("数据结构", 4.0, "A")
    t.add("大学物理", 4.0, "C")
    t.add("程序设计", 3.0, "A")

    print(t.summary())
    print("\n课程排名:")
    for name, credits, gp in t.ranking():
        print(f"  {name}: {credits}学分, 绩点{gp}")
\end{lstlisting}
\fi
\caption{跨页代码段示例（Python）}\label{lst:cross-page}
\end{longlisting}
\makeatother

行内代码可以使用 \verb|\verb| 命令，或使用高亮命令以获得更好的效果，例如：
\makeatletter
\iftongjithesis@minted
  \mintinline{python}{print("Hello")}
\else
  \lstinline[language=Python]{print("Hello")}
\fi
\makeatother
（与普通文本 \verb|print("Hello")| 对比）。

\section{分页控制}

在\LaTeX{}排版系统中，\verb|\newpage| 和 \verb|\clearpage| 命令都能强制开始新页面，但它们有着重要区别。\verb|\newpage| 仅创建新页面，而 \verb|\clearpage| 还会处理所有待处理的浮动体，确保它们在新页面前被正确放置。

当需要确保所有浮动体都已处理并开始新页面时，应优先使用 \verb|\clearpage| 而非 \verb|\newpage|，以防止浮动体累积至错误位置。
